\documentclass[10pt,a4paper]{amsart}
\usepackage[margin=19mm]{geometry}
\usepackage{amsmath,amssymb,mathtools}
\usepackage[scr=rsfs,cal=euler]{mathalfa}
\usepackage{xfrac}
\usepackage[unicode,hidelinks,bookmarks=false]{hyperref}
\newcommand{\ie}{i.e.\ }
\newcommand{\cf}{cf.\ }
\newcommand{\resp}{respectively\ }
\newcommand{\Subsection}{\subsection}
\providecommand{\nobreakdash}{\nobreak\hbox{-}}
\setlength{\emergencystretch}{2em}
\multlinegap0pt
\usepackage{bm}
\usepackage{bbm}
\usepackage{dsfont}
\usepackage{xspace}
\usepackage{cancel}
\usepackage{mathtools}
\usepackage[shortlabels]{enumitem}
\usepackage{amsthm}
\makeatletter
\@ifundefined{theorem}{\newtheorem{theorem}{Theorem}}{}
\@ifundefined{lemma}{\newtheorem{lemma}[theorem]{Lemma}}{}
\@ifundefined{proposition}{\newtheorem{proposition}[theorem]{Proposition}}{}
\makeatother
\newcommand\Ba{\mathbf B}
\newcommand\Da{\mathbf{D}}
\newcommand\Ea{\mathbf{E}}
\newcommand\Ha{\mathbf{H}}
\newcommand\alge{\mathbb{M}}
\newcommand\lu{\mathbb{L}}
\newcommand\oink{\mathcal O}
\title[Two characterizations of Bass orders via branches]{Two characterizations of Bass orders via branches}
\author{Luis Arenas-Carmona}
\date{Source: arXiv:2503.00911v2 (CC BY 4.0)}
\begin{document}
\maketitle

\noindent\emph{Source and licence.} Two characterizations of Bass orders via branches. Version arXiv:2503.00911v2 (CC BY 4.0), distributed under the Creative Commons Attribution 4.0 International licence, as recorded in the arXiv licence metadata for this exact version and shown on the abstract page https://arxiv.org/abs/2503.00911v2. Full rights notice: licenses/arxiv-2503.00911-CC-BY-4.0.txt.\par\smallskip

\noindent\textit{Editorial scope.} Counted unit: the Proposition whose source label is \texttt{p45} in the article (source0.tex lines 1354--1356, source label \texttt{p45}), delivered together with the article's own complete original proof of it (source0.tex lines 1358--1486, one \texttt{proof} environment of 129 lines). No mathematics is written, repaired or re-derived: statement and proof are the article's own text line for line. Required context retained uncounted: p26 (lines 893--908); p27 (lines 979--990); p41n (lines 1054--1079); Ea1 (lines 1057--1076); Ea2 (lines 1057--1076); l43n (lines 1303--1316); c3x (lines 1346--1351). Every definition, display and label of those lines is retained unchanged. Omitted from this package and not redistributed: 1--892, 909--978, 991--1053, 1077--1302, 1317--1345, 1352--1353, 1357--1357, 1487--2336. Nothing inside a retained excerpt is omitted or altered: every retained body is delivered exactly as the article's lines are written, so no omission receipt of any kind accompanies this package. Citations: the retained excerpts carry no citation command at all, so no bibliography entry of the article is transcribed and no reference is redistributed. Reference adaptation: none was needed and none was made; every reference inside the delivered unit resolves inside it, so no unresolved reference or citation is produced. Printed slips kept literal: no printed word, symbol, hypothesis or number of the counted statement or of its proof was altered. Layout and class: the article's own class is \texttt{amsart} with packages including babel, geometry, amsmath, amssymb, graphicx, pict2e, enumitem, tikz; the wrapper is the lane's pinned \texttt{amsart} editorial wrapper at 10pt on A4, which restates the theorem-like environments the retained text needs and adds the article's own macro definitions that the retained text needs (\texttt{\textbackslash Ba}, \texttt{\textbackslash Da}, \texttt{\textbackslash Ea}, \texttt{\textbackslash Ha}, \texttt{\textbackslash alge}, \texttt{\textbackslash lu}, \texttt{\textbackslash oink}), with extra wrapper packages (bm, bbm, dsfont, xspace, cancel, mathtools, enumitem). The delivered numbering is the unit's own. Assets: no retained span contains a figure environment, a graphics-inclusion command, a PDF-page inclusion, a table or a TikZ picture; no image file is copied into this package, the package declares no asset dependency and no image file is redistributed with it. Licence: version arXiv:2503.00911v2 of this article is distributed under the Creative Commons Attribution 4.0 International licence, as recorded in the arXiv licence metadata for this exact version and shown on the abstract page https://arxiv.org/abs/2503.00911v2. A full rights notice is delivered at licenses/arxiv-2503.00911-CC-BY-4.0.txt. Characters: every retained excerpt is pure ASCII, so the delivered engine, XeLaTeX with its default Latin Modern text font, reports no missing character.
\begin{proposition}\label{p26}
    Assume $\lu\subseteq\alge $
    is isomorphic to a ramified
    separable extension $L$ of $K$.
    Let $e$ be the edge in $\mathfrak{t}(K)$
    connecting the two maximal orders of $\alge $
    containing $\oink_{\lu}$.
    Let $\lu_L=L\lu\cong L\times L$
    be the extension of $\lu$, and let 
    $\Ha=\oink_{\lu_L}$ be its ring of integers. 
    Then the vertex $w$ of $\mathfrak{t}(L/K)$ whose distance
    $d_w$ to $\mathfrak{s}_L(\Ha)$ is minimal
    is the midpoint of the edge $e$. 
    Furthermore, $d_w=0$ if $K$ is non-dyadic, while
    $d_w>0$ for a dyadic field $K$.
\end{proposition}

\begin{proposition}\label{p27}
    Assume $\lu\subseteq\alge $
    is isomorphic to an
    inseparable extension $L$ of $K$.
    Let $\lu_L=L\lu\cong L[x]/(x^2)$.
    Let $z$ be the visual limit of the branch of
    any non-trivial order in $\lu_L$. 
    Then the vertex of $\mathfrak{t}(L/K)$
    that is closest to $z$ is the midpoint of the edge
     connecting the two maximal orders containing 
    $\oink_{\lu}$. 
\end{proposition}

\begin{proposition}\label{p41n}
Let $\Ba$ be a Bass order. Then precisely one of the following
conditions holds:
\begin{enumerate}[label=($\mathfrak{E}$\bf{\arabic*})]
    \item\label{Ea1} $\Ba$ is contained in a unique maximal 
    order $\Da$, and it has the form 
    $\Ba=\mathcal{O}_{\lu}+\pi_K^r\Da$,
    for some integer $r\geq0$, and for some two-dimensional 
    subalgebra $\lu\subseteq\alge$ that is
    isomorphic to an unramified quadratic extension $L/K$
    and satisfies $\mathcal{O}_{\lu}\subseteq\Da$.
    \item\label{Ea2} $\Ba$ is contained in two maximal orders 
    $\Da_1$ and $\Da_2$ sharing an edge, and it has the form 
    $\Ba=\mathcal{O}_{\lu}+
    \boldsymbol\pi_{\lu}^r(\Da_1\cap\Da_2)$, 
    for some algebra $\lu\subseteq\alge$ isomorphic 
    to a ramified quadratic extension $L/K$ satisfying
    $\mathcal{O}_{\lu}\subseteq(\Da_1\cap\Da_2)$
    and for some 
    uniformizer $\boldsymbol\pi_{\lu}$.
    \item \label{Ea3} $\Ba$ is an Eichler order 
    of level $d\geq 2$.
\end{enumerate} 
Furthermore, any full ideal $\Ba\subseteq\alge$ satisfying one
of these conditions is a Bass order.
\end{proposition}

\begin{lemma}\label{l43n}
Let $\mathcal{O}_\lu$ be the ring of 
integers in a subalgebra $\lu\cong L$
of $\alge=\alge $, where $L/K$ is a 
quadratic extension of local fields.
Let $\Ba$ be a full order containing 
$\mathcal{O}_\lu$. Let
$\boldsymbol{\pi}_\lu$ be a uniformizer
of $\mathcal{O}_\lu$.
Then any full suborder $\Ba'$ of $\Ba$ containing 
$\mathcal{O}_\lu$ has the form $\Ba'=
\mathcal{O}_\lu+\boldsymbol{\pi}_\lu^r\Ba$
for some $r\geq0$.
\end{lemma}

\begin{lemma}\label{c3x}
    If a maximal order $\Da$ is a leaf of the branch $\mathfrak{s}_K(\Ha)$ then the extension
    $\Da_L$ is a leaf of the branch 
    $\mathfrak{s}_L(\Ha_L)$ for any field extension
    $L/K$.   \qed 
\end{lemma}

\begin{proposition}\label{p45}
    Every Bass order is a ghost intersection of maximal orders.
\end{proposition}

\begin{proof}
We use the classification in last section 
to give a case
by case proof.  The result is well known for 
Eichler orders, for which the field $L$ 
on which the orders
are defined can be assumed to be $L=K$. Therefore,
we can assume that the Bass order $\Ba$ is 
in either of the cases \ref{Ea1} or 
\ref{Ea2} in Prop. \ref{p41n}.

Firstly, we let $\Ba$ be a Bass order of the form
$\Ba=\mathcal{O}_{\lu}+\pi_K^r\Da$,
where $L/K$ is an unramified quadratic extension,
and $\Da$ is the only maximal order containing
$\Ba$. Then the extension 
$\lu_L=L\lu\cong L\times L$
has a ring of integers 
$\oink_{\lu_L}\cong\oink_L\times\oink_L$,
which is contained precisely in the maximal orders 
in an infinite line $\mathfrak{p}$, passing through $\Da_L$, 
and whose visual limits $z$ and $\bar{z}$
satisfy $\bar{z}=\sigma(z)$,
as elements in $\mathbb{P}^1(L)$,
where $\sigma$ is the 
generator of the Galois group $\mathrm{Gal}(L/K)$.  
Let $\Da_1\subseteq\alge_L$ be a 
maximal order in $\mathfrak{p}$ at distance $r$ from $\Da$. Set 
$\Ba'=\Da_L\cap\Da_1\cap\alge $. 
This is an order in
$\alge $ that contains $\oink_{\lu}$ and 
it is contained
in $\Da$. It follows that 
$\Ba=\mathcal{O}_{\lu}+\pi_K^{r'}\Da$ for 
some integer
$r'$ by Lem. \ref{l43n}. In fact, $r'$ is the smallest integer $s$ satisfying
$\pi_K^s\Da\subseteq\Da_1$, which is equivalent
to $\Da_L^{[s]}\subseteq\Da_1$. Since 
$\mathfrak{s}_L\left(\Da_L^{[s]}\right)$ is a ball of radius $s$
around the vertex $\Da_L$, it follows that $r=r'$, 
and the result follows.

Finally, assume that $\Ba$ is a Bass order of the form 
$\Ba=\mathcal{O}_{\lu}+
\boldsymbol\pi_{\lu}^r\Ea$,
where $\Ea=\Da\cap\Da'$ is an 
Eichler order of level $1$
containing $\Ba$, the subalgebra 
$\lu\subseteq\alge $ is isomorphic
to the ramified quadratic 
extension $L$ of $K$,  and 
$\boldsymbol\pi_{\lu}\in\oink_{\lu}$
is a uniformizer. Let 
$\Ea_L=\oink_L\Ea$ be the extension,
which is an Eichler order, as it 
contains a non-trivial
idempotent. The order $\Ea_L$ is
contained precisely 
in the maximal orders between 
$\Da_L$ and $\Da_L'$ 
in the tree $\mathfrak{t}(L)$, 
as follows from Lem.
\ref{c3x}. In fact, $\Ea_L$ is contained
in exactly $3$ maximal orders in $\alge_L$,
since $\Da_L$ and $\Da_L'$ lie at distance $2$ in 
$\mathfrak{t}(L)$. The third order is
denoted $\hat{\Da}$.
It is located between $\Da_L$ and $\Da_L'$ in 
$\mathfrak{t}(L)$, and it is not defined over
$K$, i.e., is not the extension of a maximal 
order in $\alge$.
Note that we can write
$\boldsymbol\pi_{\lu}=
\pi_L\mathtt{u}$ in $\alge_L$, where 
$\mathtt{u}$ is a unit in 
the algebra $\lu_L=L\lu$,
as in the proof of Prop. \ref{p27}.
Assume first that $L/K$ is separable. Then
the matrix $\boldsymbol\pi_{\lu}$ has two 
eigenvectors in the vector 
space $L^2$, which correspond to visual 
limits $z$ and $\bar{z}$ 
of the tree $\mathfrak{t}(L)$, and the 
generator $\sigma$
of the Galois group $\mathcal{G}=
\mathrm{Gal}(L/K)$ 
satisfies $\sigma(z)=\bar{z}$. Let $\mathfrak{p}$
be the infinite line with visual limits 
$z$ and $\bar{z}$.
Let $\Da_0$ be the point of the  $\mathfrak{p}$
that is closest to $\mathfrak{t}(L/K)$.
By Prop. \ref{p26}, the vertex in 
$\mathfrak{t}(L/K)$ that is closest to $\Da_0$
is $\hat{\Da}$. Note that $\Da_0=\hat{\Da}$ 
precisely when $K$ is non-dyadic.
Consider a maximal order $\Da_1$ 
in the ray from $\hat{\Da}$ to $z$ at distance 
$r$ from $\hat{\Da}$. In the inseparable case
we use the ray joining $\hat{\Da}$ to 
the visual limit $z$ of the branch
$\mathfrak{q}$ of any non-trivial
order in $\lu_L$, as described in Prop.
\ref{p27}, but the reasoning is similar.
Note that $\oink_{\lu}$
is contained in every maximal order in the line $\mathfrak{p}$
(or the branch $\mathfrak{q}$),
and also in $\hat{\Da}$, so it must be contained in
$\Da_1$. We claim that
$\Ba=\Da_1\cap\Ea\cap\alge $. Again,
set $\Ba'=\Da_1\cap\Ea\cap\alge $.
Then $\Ba'$ is an order containing $\oink_{\lu}$
and contained in $\Ea$, so that 
$\Ba'=
\mathcal{O}_{\lu}+
\boldsymbol\pi_{\lu}^{r'}\Ea$,
for some positive integer $r'$, which is 
the smallest value of $s$ satisfying
$\boldsymbol\pi_{\lu}^s\Ea\subseteq\Da_1$.
Since $\Ea$ spans $\Ea_L$ as a $\oink_L$-module,
we can replace $\Ea$ by $\Ea_L$ and require that
$\boldsymbol\pi_{\lu}^s\Ea_L=
\pi_L^s\mathtt{u}^s\Ea_L\subseteq\Da_1$.
Since $\mathtt{u}\in\oink_{\lu}^*
\subseteq\Ea_L^*$, we can ignore it 
and just require
that $\pi_L^s\Ea_L\subseteq\Da_1$, which is 
equivalent to $s\geq r$ by the general theory. The result 
follows by reasoning as in the previous case.
\end{proof}

\end{document}
